\section{Metodología}
\subsection{El terreno y el MDP}
Usamos un recorte de unos $3\times3$~km del modelo digital de elevación (Digital
Elevation Model, DEM) Copernicus GLO-30, bajado de AWS Open Data. La
Fig.~\ref{fig:terreno} muestra la elevación y la pendiente. El MDP es:
\begin{itemize}
  \item \textbf{Estados:} el agente está en una celda de la malla $108\times108$.
  \item \textbf{Acciones:} arriba, abajo, izquierda, derecha.
  \item \textbf{Transición:} determinística, u opcionalmente estocástica (el
        resbalón de la Sec.~\ref{sec:estoc}).
  \item \textbf{Recompensa:} un costo por paso que crece con la pendiente, un
        castigo si choca contra un obstáculo, y un premio al llegar a la meta.
\end{itemize}

\subsection{Observación local}
El agente no ve el mapa completo. Ve un parche de $9\times9$ celdas de pendiente
centrado en él (81 valores) más el vector $(\Delta \text{fila}, \Delta
\text{columna})$ hacia la meta, en total 83 números (Fig.~\ref{fig:obs}). Ese
tamaño no depende del mapa. Y como terrenos parecidos dan observaciones parecidas,
la red puede generalizar. Una tabla no puede: necesitaría una entrada por cada par
(posición, meta), del orden de $10^8$ valores.

\subsection{Meta fija con inicio aleatorio}
La meta es siempre la misma celda; el inicio cambia en cada episodio. Con una sola
meta, la función de valor es un campo suave y el entrenamiento se mantiene estable.
El inicio aleatorio hace el problema no trivial: el agente aprende a llegar desde
cualquier celda, no una ruta memorizada.

\subsection{Reward shaping de progreso}
Con la meta lejos en un mapa grande, el agente casi nunca la encuentra por azar y el
premio nunca aparece. Le damos un bono por acercarse. Primero probamos el
\emph{shaping} potencial $\Phi(s)=-k\cdot d(s,\text{meta})$, y falló
(Sec.~\ref{sec:disc}). Lo cambiamos por \textbf{shaping de progreso}, $r
\mathrel{+}= k\,[\,d(s,\text{meta}) - d(s',\text{meta})\,]$, que suma cero en
cualquier ciclo y no crea la trampa del potencial.

\subsection{Arquitectura e hiperparámetros}
La red (tanto la de $Q$ como la de política) es un perceptrón de $83 \to 128 \to 128
\to 4$, unos 27~mil parámetros. Usamos $\gamma=0.99$, optimizador Adam, pérdida de
Huber y recorte de gradiente en DQN. La exploración es $\varepsilon$-greedy
decreciente. Un \emph{curriculum} hace que los inicios arranquen cerca de la meta y
se alejen con el entrenamiento. Fijamos la semilla en todos los experimentos para
que sean reproducibles.
